\section{Discrimination as an error of representation}
\label{sec:discrimination}

Nothing in \S\ref{sec:model} refers to an agent's class. We now let some agents
make a specific, minimal mistake: when they receive a message, they extend the
representation of the issue with a feature that carries no information about the
issue at all, and depends only on \emph{who sent it}. If that feature is
correlated with the emitter's class label, the agent is a
\emph{discriminating} agent, and the resulting shift of its opinion field is the
\emph{discrimination field}:
%
\begin{equation}
  \hw^{D} \;=\; \hw + D_{\kappa_r,\kappa_e},
  \label{eq:disc-field}
\end{equation}
%
where $D$ is a $2\times2$ matrix indexed by the classes of receiver and emitter.
A fraction $\fd$ of the society discriminates, drawn independently of class; the
rest use \eqref{eq:fields} unmodified.

This is a familiar failure viewed from an unusual angle: irrelevant features
degrade a classifier, and if one correlates with a protected attribute its errors
become structured along that attribute. Here the classifier is one agent among
many, and the question is what those errors do to the \emph{population}.

\paragraph{What the field does.} By \eqref{eq:modulation}, the sign of perceived
agreement decides whether trust grows or decays. A receiver that adds $+d$ to
$\hw$ reads its counterpart as more agreeable than it is and becomes more
trusting: it is \emph{tolerant}. A receiver that adds $-d$ manufactures
disagreement that was not there and becomes more distrustful: it is
\emph{intolerant}. Throughout this paper $d>0$ therefore denotes
\textbf{tolerance towards in-group and intolerance towards out-group} emitters,
and $d<0$ the reverse --- discrimination in favour of the other class.

Geometrically, the field displaces the separatrix of the learning flow to
$\hmu = \hw + D$ (Figure~\ref{fig:flows}, Appendix~\ref{app:flows}). The undisplaced flow has two
symmetric consonant basins, ``trust and agree'' and ``distrust and disagree'',
and a society whose agents fall into them polarises into two camps whose
membership has nothing to do with class. Displacing the separatrix breaks that
symmetry: a receiver meeting an out-group emitter has an enlarged
``distrust and disagree'' basin, and one meeting an in-group emitter an enlarged
``trust and agree'' basin. Every discriminating agent is thus biased towards the
same alignment of distrust with class, and once enough of them are, the whole
society can lock into it.


\paragraph{Six ways to discriminate.} The matrix $D$ can be filled in six
inequivalent ways, listed in Appendix~\ref{app:cases}: one class may favour its
own without hostility to the other, be hostile without favouritism, or both; and
either one class or both may discriminate. All the phase diagrams below use the
fully symmetric case
%
\begin{equation}
  D \;=\; d\begin{pmatrix} +1 & -1\\ -1 & +1\end{pmatrix},
  \label{eq:case6}
\end{equation}
%
in which both classes favour their own and are hostile to the other. It is the
case in which the two classes play equivalent roles, so any asymmetry that
appears in the results is spontaneous rather than imposed.

\paragraph{A note on the sign convention.} An earlier version of this model
\citep{caticha2026discrimination} tabulates $D$ with the opposite sign while
placing the discriminatory phases at $d>0$ in its text and figures; the two are
inconsistent, and the inconsistency exchanges the halves of the phase diagram.
Appendix~\ref{app:sign} works this through and shows the maps under both
readings. We use \eqref{eq:disc-field} with \eqref{eq:case6} as stated, which is
the reading that agrees with the algorithm.
